\documentclass[12pt,a4paper]{article}
\usepackage[UTF8]{ctex}
\usepackage{geometry}
\usepackage{listings}
\usepackage{xcolor}
\usepackage{hyperref}
\usepackage{booktabs}
\usepackage{graphicx}
\usepackage{amsmath}
\usepackage{float}
\usepackage{tcolorbox}
\usepackage{enumitem}
\usepackage{amssymb}

% 页面设置
\geometry{left=2.5cm,right=2.5cm,top=2.5cm,bottom=2.5cm}

% 代码高亮设置
\lstset{
    language=Python,
    basicstyle=\ttfamily\small,
    keywordstyle=\color{blue}\bfseries,
    commentstyle=\color{gray},
    stringstyle=\color{orange},
    numbers=left,
    numberstyle=\tiny\color{gray},
    stepnumber=1,
    numbersep=8pt,
    backgroundcolor=\color{gray!10},
    frame=single,
    frameround=fttt,
    breaklines=true,
    showstringspaces=false,
    tabsize=4,
    captionpos=b,
    xleftmargin=2em,
    xrightmargin=1em,
    aboveskip=1em,
    belowskip=1em
}

% 知识点框
\tcbuselibrary{skins,breakable}
\newtcolorbox{knowledgebox}[1][]{
    colback=blue!5!white,
    colframe=blue!50!black,
    fonttitle=\bfseries,
    title=#1,
    breakable
}

% 练习框
\newtcolorbox{practicebox}[1][]{
    colback=green!5!white,
    colframe=green!50!black,
    fonttitle=\bfseries,
    title=#1,
    breakable
}

% 注意框
\newtcolorbox{tipbox}[1][]{
    colback=yellow!10!white,
    colframe=orange!70!black,
    fonttitle=\bfseries,
    title=#1,
    breakable
}

% 超链接设置
\hypersetup{
    colorlinks=true,
    linkcolor=blue,
    filecolor=magenta,      
    urlcolor=cyan,
    bookmarks=true,
    pdfstartview=FitH
}

% 标题设置
\title{\textbf{Pandas入门与数据处理实战手册}\\[0.5em]\large 零基础入门版}
\author{竞赛培训组}
\date{\today}

\begin{document}

\maketitle
\tableofcontents
\newpage

%========================================
\section{准备工作：安装与导入}
%========================================

\subsection{什么是Pandas？}

\begin{knowledgebox}[知识点：Pandas是什么]
Pandas是Python的一个数据分析库，它的名字来源于"Panel Data"（面板数据）。

简单来说，Pandas就是用来处理表格数据的工具，就像Excel一样，但比Excel更强大、更灵活。

\textbf{为什么叫Pandas？}
\begin{itemize}
    \item 和熊猫（panda）没有关系
    \item 来自计量经济学的"Panel Data"概念
    \item 是Python Data Analysis的缩写
\end{itemize}
\end{knowledgebox}

\subsection{安装Pandas}

在使用Pandas之前，需要先安装它。打开命令行（CMD或PowerShell），输入以下命令：

\begin{lstlisting}[caption={安装Pandas}]
# 使用pip安装（推荐）
pip install pandas

# 如果安装了Anaconda，使用conda安装
conda install pandas
\end{lstlisting}

\begin{tipbox}[安装小贴士]
\begin{itemize}
    \item 安装时需要联网
    \item 如果安装速度慢，可以使用国内镜像源
    \item 安装完成后，关闭命令行重新打开
\end{itemize}
\end{tipbox}

\subsection{导入Pandas}

安装完成后，在Python代码中导入Pandas：

\begin{lstlisting}[caption={导入Pandas}]
# 导入pandas库，并给它起别名pd
import pandas as pd

# 同时导入numpy（经常和pandas一起使用）
import numpy as np
\end{lstlisting}

\begin{knowledgebox}[知识点：为什么要用别名？]
在Python中，给库起别名是为了：
\begin{enumerate}
    \item \textbf{简化代码}：写\texttt{pd}比写\texttt{pandas}更短
    \item \textbf{行业惯例}：大家都用\texttt{pd}，方便阅读
    \item \textbf{保持一致}：不同人写的代码风格统一
\end{enumerate}

这就好比：
\begin{itemize}
    \item 你的名字叫"张三丰"，但朋友们叫你"小张"
    \item 写代码时，\texttt{pandas}是全名，\texttt{pd}是小名
\end{itemize}
\end{knowledgebox}

%========================================
\section{认识我们的练习数据}
%========================================

\subsection{四个CSV文件介绍}

在本手册中，我们将使用四个CSV文件进行学习和练习。请确保这四个文件和你的Python代码放在同一个文件夹中。

\begin{table}[H]
\centering
\caption{练习数据文件说明}
\begin{tabular}{|l|l|l|l|}
\hline
\textbf{文件名} & \textbf{数据量} & \textbf{用途} & \